\chapter{Memoria: l'archivista bugiardo}

\begin{flushright}
	\textit{``Ogni atto del ricordare \`{e} un atto del dimenticare: rimodelliamo il passato ogni volta che lo richiamiamo, finch\'{e} non rimane pi\`{u} nulla dell'originale.''}\\[6pt]
	\textsc{\index[main]{Sacks, Oliver}Oliver Sacks}
\end{flushright}

\bigskip
\bigskip

\section{L'archivista che riscrive la storia}

Provate a immaginare di dover assumere qualcuno per custodire i vostri archivi pi\`{u} preziosi. Difficilmente scegliereste una persona che butta via documenti importanti senza avvisarvi, che ne modifica altri aggiungendo dettagli mai accaduti, che vi consegna con assoluta sicurezza fascicoli mai esistiti, e che mette sempre in primo piano le stesse tre carte lasciando tutto il resto a marcire in un cassetto. Eppure \`{e} esattamente questo il tipo di archivista che ciascuno di noi porta dentro la testa.

La memoria umana non \`{e} un disco rigido che registra fedelmente gli eventi. Assomiglia molto di pi\`{u} a un narratore compulsivo che riscrive di continuo le proprie storie per renderle pi\`{u} coerenti, pi\`{u} drammatiche o pi\`{u} lusinghiere nei confronti del protagonista. E il protagonista, naturalmente, siete sempre voi.

Questo difetto apparente \`{e} in realt\`{a} una soluzione evolutiva di rara intelligenza. Per i nostri antenati non aveva alcuna importanza ricordare con esattezza cosa avessero mangiato un certo marted\`{i}, mentre era questione di vita o di morte ricordare con assoluta nitidezza il punto esatto in cui avevano avvistato il predatore. La memoria si \`{e} evoluta per essere selettiva e teatrale proprio perch\`{e} questa selettivit\`{a} aumentava le probabilit\`{a} di sopravvivere. Il guaio \`{e} che oggi quella stessa macchina ci fa ricordare alla perfezione l'unica volta in cui un investimento \`{e} andato male, mentre cancella le dieci volte in cui era andato bene. Ci fa temere l'attentato terroristico, raro ma memorabile, molto pi\`{u} dell'incidente domestico, comune ma dimenticabile.

C'\`{e} qualcosa di affascinante in questa idea, e non \`{e} affatto nuova. Pi\`{u} di milleseicento anni fa, un uomo seduto a scrivere nel Nord Africa romano si pose le stesse domande che oggi affrontiamo con le risonanze magnetiche. Agostino d'Ippona, nel decimo libro delle sue \textit{Confessioni}, descrisse la memoria come ``i vasti palazzi della memoria'', un luogo sterminato in cui custodiamo non semplici registrazioni, ma il nostro stesso essere.\footnote{Agostino d'Ippona (398 d.C. circa). \textit{Confessioni}, Libro X. Numerose edizioni.} La sua intuizione pi\`{u} sorprendente fu quella di non considerare la memoria un magazzino del passato, bens\`{i} ``il presente del passato'': il passato, per Agostino, non esiste pi\`{u} da nessuna parte se non nel modo in cui lo facciamo rivivere ora, in questo istante. Quando ricordiamo, non apriamo un cassetto polveroso. Costruiamo qualcosa di nuovo, qui e adesso, con i materiali che abbiamo a disposizione oggi.

Sedici secoli dopo, la neuroscienza non ha fatto che dare ragione a quel vescovo africano. Ma per capire fino in fondo quanto sia inaffidabile la nostra memoria, dobbiamo scendere nei dettagli di come funziona davvero questo sistema bizzarro.

\section{L'illusione del cassetto: come si ricostruisce un ricordo}

Cosa succede esattamente, dentro il vostro cranio, nel momento in cui rievocate qualcosa? La risposta capovolge tutto ci\`{o} che istintivamente crediamo. Quando richiamate un ricordo non state estraendo un documento intatto da un archivio: lo state letteralmente ricostruendo da zero, come un archeologo che, trovati pochi cocci di ceramica, deve immaginare la forma dell'intero vaso.

Il cervello afferra alcuni punti di ancoraggio dell'esperienza originale e li ricollega tra loro usando le conoscenze, le emozioni e le aspettative che avete \textit{ora}, in questo preciso momento. \`{E} un'improvvisazione costruita su frammenti dispersi, in cui i vuoti vengono riempiti con ci\`{o} che oggi vi sembra plausibile, non necessariamente con ci\`{o} che era vero allora. Il dettaglio pi\`{u} sconcertante \`{e} che il ricordo non \`{e} nemmeno conservato in un unico luogo. Risulta frammentato in pezzi sparsi per tutto il cervello: i colori in un'area, i suoni in un'altra, le emozioni in un'altra ancora, gli odori in una regione completamente diversa.\footnote{Moscovitch, M., e altri (2016). Episodic memory and beyond: The hippocampus and neocortex in transformation. \textit{Annual Review of Psychology}, 67.} Quando ricordate, il cervello deve recuperare tutti questi frammenti da magazzini lontani e rimontarli al volo, come un regista che assembla un film in tempo reale partendo da bobine custodite in depositi diversi.

\begin{figure}[htbp]
	\centering
	\includegraphics[width=0.9\textwidth]{images/infografica_cap3_1.png}
	\captionof{figure}{La frammentazione della memoria: come un singolo ricordo viene scomposto e distribuito in diverse aree cerebrali.}
	\label{fig:memoria_frammentata}
\end{figure}

Questa frammentazione spiega perch\'{e} a volte ricordate alla perfezione l'odore di un momento ma non il volto della persona che vi stava accanto, oppure perch\'{e} un profumo improvviso pu\`{o} riportarvi alla mente un intero episodio che credevate sepolto per sempre. Ogni frammento pu\`{o} essere richiamato in modo indipendente, e alcuni resistono all'oblio molto pi\`{u} di altri. C'\`{e} perfino un curioso rovesciamento nell'ordine: mentre nella percezione originale i dettagli visivi precedono il significato, durante il recupero il cervello ricostruisce prima il ``cosa'', il significato, e solo dopo il ``come appariva'', i dettagli percettivi.\footnote{Linde-Domingo, J., e altri (2019). Evidence that neural information flow is reversed between object perception and object reconstruction from memory. \textit{Nature Communications}, 10(1).} \`{E} come se il vostro cervello dicesse: ``So che era tua nonna, adesso lasciami ricostruire com'era vestita.''

Tutto questo avviene in modo cos\`{i} automatico e convincente che restiamo assolutamente certi dell'accuratezza dei nostri ricordi, persino quando sono palesemente falsi. Il cervello compila le lacune, integra informazioni sopraggiunte in seguito, elimina le incongruenze, e lo fa con tale fluidit\`{a} che il risultato ci appare perfettamente genuino. Non percepiamo nessuna manipolazione perch\'{e} l'intera operazione si svolge molto al di sotto della soglia della consapevolezza.

Prendiamo un esempio che riguarda tutti. Pensate al vostro primo giorno di scuola. Probabilmente conservate ricordi piuttosto vividi: come eravate vestiti, cosa avete provato, qualche particolare dell'aula o dei compagni. Eppure molto di ci\`{o} che ``ricordate'' con ogni probabilit\`{a} non \`{e} mai accaduto, o almeno non esattamente cos\`{i}. Quel ricordo \`{e} un collage costruito con frammenti reali dell'esperienza originale, mescolati a elementi presi da altri primi giorni di scuola, a storie raccontate dai genitori, a fotografie viste in seguito, persino a scene di film o libri assorbite inconsapevolmente. Il cervello ha preso tutti questi pezzi e li ha assemblati in una narrazione coerente, aggiungendo colori pi\`{u} accesi, emozioni pi\`{u} intense, dettagli pi\`{u} precisi di quelli che un bambino di cinque o sei anni avrebbe potuto cogliere. Ha creato, in sostanza, una versione romanzata della vostra esperienza reale. E la cosa pi\`{u} straordinaria \`{e} che questo ricordo ricostruito pu\`{o} apparirvi pi\`{u} reale e dettagliato dell'esperienza originale.

\section{Il narratore necessario: perch\'{e} dimenticare \`{e} un'arte}

A questo punto sorge spontanea una domanda: perch\'{e} mai la memoria dovrebbe funzionare in un modo tanto inefficiente? La risposta, come quasi sempre, sta nell'evoluzione. Una memoria perfetta non \`{e} soltanto impossibile dal punto di vista energetico, perch\'{e} richiederebbe un cervello mostruoso per immagazzinare ogni dettaglio di ogni esperienza, ma sarebbe addirittura controproducente.

Immaginate di ricordare ogni cosa con accuratezza assoluta: ogni conversazione tediosa, ogni momento di noia, ogni piccolo fastidio fisico, ogni particolare irrilevante di ogni singola giornata. Restereste paralizzati sotto il peso dell'informazione, incapaci di distinguere ci\`{o} che conta da ci\`{o} che non conta. La memoria selettiva, al contrario, vi permette di fare la cosa evolutivamente pi\`{u} importante: imparare dalle esperienze passate per orientarvi meglio nel futuro. Non avete bisogno di ricordare cosa avete mangiato ogni giorno della vita, vi basta ricordare quali cibi vi hanno fatto star male. La memoria \`{e} stata ottimizzata per essere un sistema di apprendimento, non un sistema di archiviazione. E per imparare bene, conta pi\`{u} avere ricordi utili che ricordi accurati.

Qui entra in gioco un pensatore che, alla fine dell'Ottocento, comprese tutto questo prima ancora che esistessero gli strumenti per dimostrarlo. Henri Bergson, nel suo \textit{Materia e memoria}, distinse due forme radicalmente diverse di ricordo.\footnote{Bergson, H. (1896). \textit{Mati\`{e}re et m\'{e}moire}. F\'{e}lix Alcan, Parigi.} Da una parte c'\`{e} la memoria abitudine, quella del corpo: \`{e} il sapere andare in bicicletta, il recitare una poesia imparata a forza di ripetizioni, un meccanismo automatico inscritto nei gesti. Dall'altra c'\`{e} la memoria ricordo, quella viva e spirituale, che fa rivivere un singolo momento irripetibile del passato. Per Bergson la memoria non era affatto un deposito immobile di immagini: era una forza creativa che filtra e seleziona il passato in funzione di ci\`{o} che serve all'azione presente. Il cervello, sosteneva, non serve tanto a conservare i ricordi quanto a tenerne fuori la maggior parte, lasciando affiorare solo quelli utili al momento. L'oblio, in altre parole, non \`{e} il nemico della memoria. \`{E} la sua condizione di funzionamento.

La scienza degli ultimi anni ha confermato questa visione in modo quasi letterale, scoprendo qualcosa di radicale: l'oblio non \`{e} la semplice assenza del ricordare, non \`{e} il lato oscuro della memoria. \`{E} una funzione cognitiva autonoma, dotata di una propria identit\`{a} e di un proprio scopo evolutivo.\footnote{Richards, B. A., e Frankland, P. W. (2017). The persistence and transience of memory. \textit{Neuron}, 94(6).} Pensate al cameriere che prende il vostro ordine: se ricordasse tutti i caff\`{e} serviti negli ultimi dieci anni non riuscirebbe a portarvi il vostro. Ogni mattina, quando parcheggiate l'auto, dovete dimenticare attivamente dove l'avevate lasciata ieri, altrimenti il cervello vi restituirebbe un mosaico confuso di centinaia di parcheggi sovrapposti. L'oblio non \`{e} un guasto del sistema: \`{e} il sistema che si aggiorna di continuo per restare funzionale.

Il cervello, dunque, lavora attivamente per cancellare, non si limita a dimenticare per pigrizia. Le curve dell'oblio descritte da Ebbinghaus gi\`{a} nel 1885 rivelano un andamento preciso: si dimentica moltissimo all'inizio, con un crollo verticale nelle prime ore, poi la curva si appiattisce lentamente.\footnote{Ebbinghaus, H. (1885). \textit{\"Uber das Ged\"achtnis}. Duncker und Humblot, Lipsia.} \`{E} come se il cervello facesse una prima scrematura brutale, ``questo di sicuro non mi serve'', per poi diventare pi\`{u} prudente con ci\`{o} che resta. Ma questo crollo iniziale pu\`{o} essere contrastato. La chiave si chiama \index[cognitive]{generation effect}\textbf{effetto generazione}: ci\`{o} che producete attivamente, ripetendo ad alta voce, spiegando a qualcun altro, riscrivendo con parole vostre, si fissa molto meglio di ci\`{o} che leggete passivamente.\footnote{Slamecka, N. J., e Graf, P. (1978). The generation effect: Delineation of a phenomenon. \textit{Journal of Experimental Psychology: Human Learning and Memory}, 4(6).} Non basta sottolineare o rileggere: bisogna costringere il cervello a ricostruire l'informazione. E se lo fate a intervalli distanziati nel tempo, quel ricordo comincia ad aggrapparsi ad altre conoscenze e diventa parte stabile del vostro bagaglio. Nessuno ricorda quando ha imparato che Parigi \`{e} la capitale della Francia: \`{e} semplicemente diventato parte di noi.


Ed eccoci all'aspetto pi\`{u} vertiginoso di questo sistema: la sua capacit\`{a} di automigliorarsi. Ogni volta che richiamate un ricordo non vi limitate a ricostruirlo, lo modificate leggermente per renderlo pi\`{u} coerente con ci\`{o} che sapete adesso o pi\`{u} utile alla situazione presente. Se avete scoperto qualcosa di nuovo su una persona, questa informazione pu\`{o} retrocontaminare i ricordi che avevate di lei. Se avete cambiato opinione su un evento, i vostri ricordi di quell'evento possono lentamente modificarsi per allinearsi alla nuova prospettiva. Il processo \`{e} cos\`{i} sottile e graduale che non ve ne accorgete: credete di ricordare sempre la stessa cosa, mentre i vostri ricordi stanno silenziosamente evolvendo per rimanere rilevanti. Era esattamente questo che intuiva Bergson, ed esattamente questo che Agostino chiamava il presente del passato. Non possedete un passato come si possiede un oggetto: lo rifate ogni volta che lo pensate.

\section{Il bisogno di una trama: memoria, identit\`{a}, popoli}

A questo punto dobbiamo affrontare la forza pi\`{u} potente che governa l'intero meccanismo, quella che decide cosa conservare e cosa lasciar svanire. Non \`{e} la logica. \`{E} l'emozione, e dietro l'emozione un bisogno ancora pi\`{u} profondo: il bisogno di significato.

Non ricordiamo gli eventi in modo neutro: li ricordiamo attraverso il filtro di ci\`{o} che abbiamo provato. E pi\`{u} intense erano le emozioni, pi\`{u} vivido e duraturo sar\`{a} il ricordo. Anche questo ha perfettamente senso dal punto di vista evolutivo, perch\'{e} gli eventi emotivamente intensi, gli attacchi dei predatori, i conflitti sociali, gli incontri con potenziali partner, erano proprio quelli in cui bisognava imparare in fretta e ricordare bene. Il problema \`{e} che il cervello emotivo non distingue tra emozioni utili ed emozioni problematiche: ricorda con la stessa intensit\`{a} il primo bacio e un incidente stradale, il giorno della laurea e quello di un'umiliazione pubblica. Cos\`{i} i nostri ricordi pi\`{u} vividi non sono quasi mai i pi\`{u} importanti o rappresentativi della nostra vita, ma semplicemente i pi\`{u} carichi di emozione, mentre i giorni tranquilli scivolano via senza lasciare traccia.

Su questo punto Friedrich Nietzsche scrisse pagine di una durezza memorabile. Nella \textit{Genealogia della morale} si pose una domanda quasi feroce: come ha fatto un animale dominato dall'istinto e dall'oblio come l'uomo a costruirsi una memoria?\footnote{Nietzsche, F. (1887). \textit{Zur Genealogie der Moral}. C. G. Naumann, Lipsia.} La sua risposta \`{e} disturbante e illuminante insieme. L'uomo, per Nietzsche, \`{e} l'animale che pu\`{o} fare promesse, l'unico capace di vincolare il proprio futuro impegnandosi nel presente. Ma una facolt\`{a} cos\`{i} contro natura ha richiesto un prezzo spaventoso: ``si imprime qualcosa a fuoco, perch\'{e} rimanga nella memoria''. Solo ci\`{o} che non smette di far male, sosteneva, resta davvero impresso. Il dolore \`{e} stato il pi\`{u} potente strumento mnemonico della storia umana. \`{E} esattamente ci\`{o} che oggi la neurobiologia descrive con linguaggio asettico, parlando di cortisolo che amplifica la formazione dei ricordi: ricordiamo le ferite molto meglio delle carezze, perch\'{e} il cervello, proprio come immaginava Nietzsche, scrive le sue lezioni pi\`{u} importanti con il ferro rovente.

Ma c'\`{e} una dimensione ulteriore, che trasforma la memoria individuale in qualcosa di collettivo. Non ricordiamo soltanto le nostre esperienze: assorbiamo costantemente i ricordi degli altri, spesso senza accorgercene. Quando qualcuno vi racconta un episodio del passato, il cervello non si limita ad archiviarlo come storia di seconda mano: lo visualizza, lo rivive, e finisce spesso per integrarlo nei vostri ricordi personali. \`{E} cos\`{i} che a volte siamo convinti di ricordare scene a cui non abbiamo mai assistito. Il fenomeno \`{e} fortissimo per l'infanzia, dove molti dei ricordi pi\`{u} antichi sono in realt\`{a} un impasto di esperienze vissute e racconti familiari ascoltati negli anni successivi. La vostra memoria personale \`{e}, in larga parte, una memoria collaborativa, coscritta da tutte le persone significative che avete incontrato. E anche questo, nei piccoli gruppi in cui si sono evoluti i nostri antenati, era un vantaggio di sopravvivenza: se un membro della trib\`{u} aveva incontrato un pericolo, conveniva che tutti gli altri lo ``ricordassero'' come se l'avessero vissuto.

Questo meccanismo ancestrale si \`{e} evoluto in qualcosa di molto pi\`{u} complesso e insidioso: la memoria collettiva delle societ\`{a} moderne. Ed \`{e} qui che emerge una distinzione fondamentale, che lo storico \textit{Alessandro Barbero} non si stanca di ribadire: storia e memoria non sono la stessa cosa, anzi, spesso si combattono apertamente.\footnote{Barbero, A. (2022). \textit{Come scoppiano le guerre}. Laterza, Roma e Bari.} La storia cerca faticosamente di ricostruire cosa \`{e} successo davvero, setacciando documenti, incrociando fonti, verificando date. La memoria, invece, se ne infischia della verit\`{a} fattuale: vuole significato, identit\`{a}, consolazione.

Prendiamo il Medioevo, i famosi ``secoli bui''. La storia documentata rivela un periodo di straordinaria vitalit\`{a} intellettuale e tecnologica: l'invenzione degli occhiali consent\`{i} agli studiosi di continuare a leggere in et\`{a} avanzata, l'aratro pesante rivoluzion\`{o} l'agricoltura, i mulini ad acqua fornirono energia meccanica su scala mai vista. Le prime universit\`{a} europee, Bologna nel 1088, Oxford intorno al 1096, Parigi verso il 1150, nascono proprio in questi presunti secoli di ignoranza. I monasteri funzionavano come centri di ricerca, dove i monaci sperimentavano tecniche agricole, sviluppavano la notazione musicale moderna, perfezionavano la distillazione.\footnote{Gies, F., e Gies, J. (1994). \textit{Cathedral, Forge, and Waterwheel: Technology and Invention in the Middle Ages}. HarperCollins, New York.} La memoria collettiva, per\`{o}, ha costruito tutt'altra narrazione: mille anni di tenebre e superstizione, un buco nero tra la gloria di Roma e lo splendore del Rinascimento. Non \`{e} che questa distorsione sia nata per caso. Ha svolto una funzione precisa: permettere prima al Rinascimento di presentarsi come ``rinascita'', poi all'Illuminismo di celebrare il trionfo della Ragione sulla superstizione. La memoria collettiva aveva bisogno di un antagonista per rendere eroica la propria storia, e il Medioevo, abbastanza lontano da non potersi difendere e abbastanza vicino da servire da monito, fu perfetto per quel ruolo.

\section{Gli inganni quotidiani e l'io che si racconta}

Il cervello individuale compie la stessa identica operazione su scala personale. La vostra storia, ci\`{o} che \`{e} accaduto davvero, potrebbe forse essere ricostruita attraverso email, fotografie, testimonianze altrui. Ma la vostra memoria di quegli stessi eventi \`{e} un'altra cosa: \`{e} il romanzo che raccontate a voi stessi per dare senso a chi siete. E quando storia e memoria entrano in conflitto, indovinate chi vince? La memoria, sempre. Perch\'{e} il cervello non \`{e} uno storico in cerca di verit\`{a}, \`{e} un narratore in cerca di sopravvivenza psicologica. Barbero lo dice senza giri di parole: i popoli ricordano ci\`{o} che conviene loro ricordare. Lo stesso vale per i singoli. Un popolo che ricordasse con precisione ogni massacro commesso e subito sarebbe paralizzato dal peso del passato; un individuo che ricordasse oggettivamente ogni fallimento e ogni umiliazione non riuscirebbe ad alzarsi dal letto.

\`{E} qui che entra in scena uno dei pensatori pi\`{u} acuti del Novecento, il filosofo francese \textit{Paul Ricoeur}. Nella sua opera \textit{S\'{e} come un altro} svilupp\`{o} il concetto di identit\`{a} narrativa: noi non siamo una sostanza fissa, una pietra che resta identica a se stessa nel tempo.\footnote{Ricoeur, P. (1990). \textit{Soi-m\^{e}me comme un autre}. \'{E}ditions du Seuil, Parigi.} Noi siamo la storia che raccontiamo di noi stessi. La nostra identit\`{a} non \`{e} data una volta per tutte: \`{e} una trama che intrecciamo di continuo, una narrazione che annoda insieme gli episodi sparsi della nostra esistenza per farne una vita sola, riconoscibile come ``mia''. Ricoeur coglie cos\`{i} il punto decisivo: se la memoria riscrive di continuo, allora l'identit\`{a} \`{e} sempre, in qualche misura, una finzione. Ma attenzione, non una bugia. Una finzione nel senso nobile del termine: una composizione, un montaggio narrativo che ci permette di rimanere noi stessi pur cambiando in continuazione. L'archivista bugiardo di cui parliamo \`{e}, in fondo, l'autore segreto della vostra autobiografia.

Questo bisogno di una trama coerente alimenta una delle caratteristiche pi\`{u} insidiose della memoria, quello che potremmo chiamare il \index[cognitive]{bias di coerenza}\textbf{bias di coerenza}. Non ci basta che i ricordi abbiano un senso emotivo: devono incastrarsi alla perfezione nella narrazione complessiva della nostra vita. \`{E} come se il nostro archivista interno non tollerasse contraddizioni nella trama. Se oggi vi considerate persone da sempre coraggiose, i ricordi dei vostri episodi di codardia tendono a svanire o a essere reinterpretati come ``prudenza''. Se vi vedete come individui di buon gusto, i ricordi dei vostri errori estetici diventano sfocati o vengono giustificati con un ``allora andava di moda''. Il risultato \`{e} che gran parte di noi conserva ricordi pi\`{u} coerenti, razionali e positivi di quanto fossero le esperienze originali. Pu\`{o} sembrare un difetto, ma dal punto di vista psicologico \`{e} una risorsa preziosa: il bias di coerenza ci consente di mantenere un'identit\`{a} stabile e un'autostima vivibile, di fronte alla naturale incoerenza della vita reale. \`{E} la finzione coerente di Ricoeur che lavora silenziosamente, fotogramma per fotogramma.

I bias mnemonici, del resto, non sono soltanto stampelle per proteggere l'ego: sono strumenti cognitivi fondamentali senza i quali non potremmo nemmeno pensare. Senza di essi non sapremmo neppure riconoscere un cane per strada. Come fate a sapere che quel bassotto appartiene alla stessa categoria del pastore tedesco? Pixel per pixel sono diversissimi, eppure il cervello li classifica entrambi come ``cani'' ignorando deliberatamente i dettagli che li distinguono. Questa categorizzazione richiede di dimenticare: dovete non vedere le differenze per cogliere la somiglianza. I rarissimi individui con una memoria troppo perfetta, come il celebre mnemonista studiato da Lurija o il Funes immaginato da Borges, non riescono a formare categorie.\footnote{Lurija, A. R. (1968). \textit{Malen'kaja knizhka o bol'shoj pamjati}. Mosca.} Sommersi dai dettagli, ricordano ogni singola foglia di ogni albero ma non arrivano mai al concetto astratto di ``albero''. Senza oblio non c'\`{e} astrazione; senza astrazione non c'\`{e} ragionamento. La vostra capacit\`{a} di pensare dipende, alla lettera, dalla vostra capacit\`{a} di dimenticare i dettagli irrilevanti per trattenere il pattern. L'intelligenza non sta nel ricordare tutto: sta nel dimenticare in modo intelligente.

\section{Il terapeuta pigro: quando dimenticare \`{e} compassione}

C'\`{e} un aspetto della memoria che raramente consideriamo: la sua funzione protettiva. Il vostro archivista non \`{e} solo disordinato, \`{e} anche uno psicologo che lavora senza sosta per preservare la vostra salute mentale. Dimenticare, da questo punto di vista, non \`{e} un guasto: \`{e} un atto di clemenza.

Il fenomeno dell'amnesia dissociativa dopo eventi traumatici non \`{e} il cervello che si rompe, \`{e} il cervello che vi protegge da ricordi capaci di paralizzarvi.\footnote{van der Kolk, B. A. (2014). \textit{The Body Keeps the Score}. Viking, New York.} Pensateci: se conservaste con perfetta nitidezza ogni istante di dolore fisico intenso, probabilmente non osereste pi\`{u} fare nulla. Le donne che hanno partorito spesso raccontano di ricordare che faceva male, ma di non riuscire a richiamare la sensazione esatta del dolore. Se la ricordassero alla perfezione, la specie umana si sarebbe forse gi\`{a} estinta.\footnote{Niven, C. A. (1988). Labour pain: Long-term recall and consequences. \textit{Journal of Reproductive and Infant Psychology}, 6(2).} Questo psicologo interno svolge anche un lavoro pi\`{u} sottile: attenua gradualmente l'intensit\`{a} emotiva dei ricordi negativi mentre preserva, o perfino amplifica, quella dei ricordi positivi. \`{E} il cosiddetto \index[main]{memoria selettiva}\textit{memoria selettiva}: le emozioni spiacevoli sbiadiscono pi\`{u} in fretta di quelle piacevoli.\footnote{Walker, W. R., Skowronski, J. J., e Thompson, C. P. (2003). Life is pleasant, and memory helps to keep it that way. \textit{Review of General Psychology}, 7(2).} Dopo una rottura devastante, all'inizio ricordate ogni dettaglio doloroso. Col tempo, per\`{o}, la memoria comincia il suo editing terapeutico: sfuma i momenti peggiori, evidenzia quelli buoni, e trasforma lentamente un trauma in una storia di crescita. Non \`{e} negazione, \`{e} sopravvivenza.

Esiste un teatro ancora pi\`{u} estremo in cui questo psicologo mostra la sua maestria: il mondo dei sogni. Ogni notte, per circa due ore complessive, viviamo avventure elaborate e bizzarre, e poi, nel giro di pochi minuti dal risveglio, il 95\%  di quei ricordi svanisce come nebbia al sole.\footnote{Hobson, J. A. (2009). REM sleep and dreaming: Towards a theory of protoconsciousness. \textit{Nature Reviews Neuroscience}, 10(11).} Non \`{e} negligenza, \`{e} progetto. Durante il sonno REM, quando i sogni sono pi\`{u} vividi, i livelli di noradrenalina, il neurotrasmettitore chiave per fissare i ricordi, crollano a zero: \`{e} l'unico momento della giornata in cui questo accade. Il cervello spegne letteralmente il registratore proprio mentre proietta i suoi film pi\`{u} surreali.\footnote{Pace-Schott, E. F., e altri (2015). Physiological feelings in dreams: Altered state or altered memory? \textit{Consciousness and Cognition}, 38.} Immaginate l'alternativa: se ricordaste ogni sogno con la chiarezza delle esperienze diurne, dopo qualche anno avreste un archivio mentale dove realt\`{a} e fantasia diventerebbero indistinguibili. Avreste ``ricordi'' di voli, di conversazioni con i morti, di trasformazioni impossibili, tutti vividi quanto quelli reali. \`{E} lo stesso principio del \textit{monitoraggio della realtà} che il cervello applica di continuo: deve separare ci\`{o} che \`{e} accaduto davvero da ci\`{o} che avete immaginato, temuto o desiderato.\footnote{Johnson, M. K., e Raye, C. L. (1981). Reality monitoring. \textit{Psychological Review}, 88(1).}

Ma tra protezione e patologia il confine \`{e} sottile. Quando questo meccanismo diventa eccessivo, potrebbe non essere pi\`{u} il cervello che vi protegge, ma il cervello che inizia a perdere colpi. Il bias della positivit\`{a} negli anziani, quella tendenza a vedere il mondo attraverso lenti rosa, potrebbe essere meno saggezza accumulata e pi\`{u} un segnale d'allarme. \`{E} romantico pensare che sia l'esperienza di vita a renderli sereni, ma la realt\`{a} neurologica suggerisce qualcosa di pi\`{u} inquietante: con l'et\`{a}, le aree cerebrali deputate al riconoscimento delle emozioni negative cominciano a deteriorarsi. Non \`{e} che gli anziani scelgano di vedere il mondo pi\`{u} buono, \`{e} che il loro cervello sta perdendo la capacit\`{a} di decodificare i segnali di pericolo. La differenza con la salutare \textit{memoria selettiva} sta nel controllo: un cervello sano sa ancora vedere la negativit\`{a} quando serve, e poi sceglie come gestirla; un cervello che scivola verso la demenza non ha pi\`{u} questa scelta, vede positivit\`{a} perch\'{e} \`{e} l'unica cosa che riesce ancora a decifrare, come una vecchia radio che riceve una sola stazione.

\section{Il tempo elastico e il montaggio della nostalgia}

Avete mai notato come le estati dell'infanzia sembrassero infinite, mentre gli anni adulti volano via in un soffio? Non \`{e} solo nostalgia poetica, \`{e} matematica neurologica. E rivela qualcosa di sorprendente: la memoria non edita soltanto i contenuti, distorce la percezione stessa del tempo.

Il cervello misura il tempo in base alla quantit\`{a} di ricordi nuovi che riesce a formare. A otto anni quasi tutto \`{e} novit\`{a}, e ogni giornata deposita strati densi di esperienze inedite. A quarant'anni la routine ha preso il sopravvento, e settimane intere si comprimono in un unico ricordo generico, ``sono andato al lavoro''.\footnote{Wittmann, M., e Lehnhoff, S. (2005). Age effects in perception of time. \textit{Psychological Reports}, 97(3).} \`{E} il cosiddetto paradosso della vacanza: due settimane in un luogo nuovo, piene di esperienze inedite, sembrano durare un'eternit\`{a} mentre le vivete, ma due settimane di routine scivolano via senza lasciare traccia.\footnote{Eagleman, D. M. (2008). Human time perception and its illusions. \textit{Current Opinion in Neurobiology}, 18(2).} La differenza sta nella densit\`{a} di marcatori temporali unici che il cervello pu\`{o} usare come punti di riferimento.

\begin{figure}[htbp]
	\centering
	\includegraphics[width=0.8\textwidth]{images/Perché_gli_anni_volano_.png}
	\captionof{figure}{La compressione temporale con l'et\`{a}: man mano che diminuisce la proporzione di esperienze nuove, gli anni sembrano accelerare. Un anno a 10 anni rappresenta il 10\% della vita vissuta; a 60 anni, appena l'1,6\%.}
	\label{fig:tempo_soggettivo}
\end{figure}

\begin{comment}
	contenuto
\noindent\fbox{\begin{minipage}{0.95\textwidth}
		\textbf{Il paradosso della vacanza: perch\'{e} due settimane sembrano eterne (e poi svaniscono)}
		
		Durante la vacanza: ``Quando finisce? Sono passati solo tre giorni e sembra un'eternit\`{a}.''
		
		Sei mesi dopo: ``Non posso credere che sia gi\`{a} agosto. Dove sono finiti questi mesi?''
		
		Non \`{e} contraddittorio: sono due orologi cerebrali diversi al lavoro.
		
		\textbf{Tempo prospettico} (mentre lo vivete): si basa sull'attenzione e sull'aspettativa. In vacanza ogni giorno \`{e} pieno di novit\`{a}, ogni momento chiede attenzione, e il tempo sembra rallentare.
		
		\textbf{Tempo retrospettivo} (quando lo ricordate): si basa sulla quantit\`{a} di ricordi immagazzinati. Quelle due settimane hanno prodotto una densit\`{a} enorme di memorie uniche, e a posteriori sembrano mesi.
		
		Al contrario, due settimane di routine volano mentre le vivete, perch\'{e} siete in pilota automatico, ma scompaiono dalla memoria, perch\'{e} non lasciano ricordi distintivi.
		
		\textbf{La strategia:} volete che l'anno non vi scivoli via? Riempitelo di novit\`{a}, anche minime. Cambiate percorso per andare al lavoro. Provate un ristorante nuovo. Imparate tre parole in una lingua straniera. Non serve partire per Bali: serve rompere la routine quel tanto che basta a creare marcatori temporali unici.
\end{minipage}}

\medskip
...
\end{comment}
Il cervello usa dunque due orologi completamente diversi: uno per vivere il presente, fondato sull'attenzione, e uno per ricordare il passato, fondato sulla quantit\`{a} di informazione immagazzinata.\footnote{Block, R. A., e Zakay, D. (1997). Prospective and retrospective duration judgments: A meta-analytic review. \textit{Psychonomic Bulletin and Review}, 4(2).} Ecco perch\'{e} il detto ``il tempo vola quando ti diverti'' \`{e} vero solo a met\`{a}: vola nel presente, ma si dilata nel ricordo. Lo stesso principio spiega perch\'{e} i momenti di pericolo estremo sembrano svolgersi al rallentatore. Durante un incidente d'auto il cervello passa in modalit\`{a} registrazione ad alta definizione, l'amigdala prende il controllo, ogni millisecondo viene impresso con precisione cinematografica. Non \`{e} che il tempo rallenti davvero: \`{e} che state formando ricordi cos\`{i} densi che, a posteriori, quei tre secondi sembreranno durare minuti.\footnote{Stetson, C., Fiesta, M. P., e Eagleman, D. M. (2007). Does time really slow down during a frightening event? \textit{PLoS One}, 2(12).}

A questa distorsione del tempo si aggiunge un montaggio ancora pi\`{u} sottile, quello della nostalgia. Il cervello ricorda con precisione cristallina gli inizi e le fini delle esperienze, mentre lascia svanire tutto quello che sta in mezzo.\footnote{Murdock, B. B. (1962). The serial position effect of free recall. \textit{Journal of Experimental Psychology}, 64(5).} \`{E} come se il vostro montatore interno conservasse soltanto la prima e l'ultima scena di ogni film, scartando i tempi morti. Ecco perch\'{e} le estati dell'infanzia sembravano perfette: ricordate l'eccitazione dell'ultimo giorno di scuola e la malinconia del rientro a settembre, mentre i quaranta giorni di noia nel mezzo, le ore passate a fissare il soffitto, i pomeriggi troppo caldi, sono evaporati come se non fossero mai esistiti.\footnote{Fredrickson, B. L., e Kahneman, D. (1993). Duration neglect in retrospective evaluations of affective episodes. \textit{Journal of Personality and Social Psychology}, 65(1).} La memoria fa quello che farebbe un bravo regista: taglia le parti noiose e tiene solo i momenti che contano. Lo stesso vale per le relazioni passate. Vi siete mai chiesti perch\'{e} a volte rimpiangiamo persone che, razionalmente, sappiamo essere state tossiche? Il cervello conserva il primo bacio e l'ultima carezza, mentre i mesi di incomprensioni quotidiane si dissolvono. Non \`{e} masochismo, \`{e} montaggio.

Il meccanismo ci permette di guardare indietro alla nostra vita e vedere una narrazione coerente e significativa, invece di ci\`{o} che spesso era, una sequenza per lo pi\`{u} ordinaria intervallata da rare punte di intensit\`{a}. Il nostro archivista bugiardo non si limita a proteggerci dal dolore: ci regala anche una versione cinematografica del nostro passato, con tutti i tagli giusti per farlo sembrare un capolavoro anzich\'{e} un documentario amatoriale. Ed \`{e} la stessa identica operazione che Ricoeur attribuiva all'identit\`{a} narrativa: trasformare il caos degli eventi in una trama che si possa abitare.

\section{Il montatore invisibile e l'animale dell'oblio}

Tutto questo ci conduce a una conclusione tanto affascinante quanto inquietante: la vostra autobiografia mentale non \`{e} un documentario, \`{e} un film d'autore in cui voi siete soltanto gli attori protagonisti, mentre il cervello fa da regista, montatore e direttore della fotografia. Il vostro cervello non vuole la verit\`{a} dei fatti: vuole una storia che vi faccia alzare dal letto la mattina. Una narrazione utile alla sopravvivenza, psicologicamente sostenibile, socialmente accettabile e, soprattutto, narrativamente soddisfacente. L'accuratezza storica? \`{E} sacrificabile, come le scene tagliate di un film per rispettare i tempi.

Questo non significa che tutti i vostri ricordi siano inventati di sana pianta. Significa che sono tutti parzialmente rimontati, selettivamente ritoccati nei colori, creativamente ridoppiati. Il materiale originale c'\`{e}, ma \`{e} passato dalla sala di montaggio, dove ha ricevuto effetti speciali emotivi, filtri nostalgici e una colonna sonora di giustificazioni aggiunte in seguito. \`{E} come guardare un vecchio film restaurato: i fotogrammi originali esistono ancora, ma qualcuno ha ravvivato i colori sbiaditi, eliminato i graffi, sincronizzato meglio l'audio, tagliato le scene che rallentavano il ritmo. Il film che vedete \`{e} ``quello originale'' e insieme profondamente trasformato. Giudicato come storico, il nostro archivista fa un lavoro pessimo. Ma come narratore, come psicologo protettivo, come consulente per la sopravvivenza esistenziale, \`{e} semplicemente insuperabile. 


S\`{i}, il nostro archivista \`{e} un bugiardo, ma \`{e} un bugiardo necessario. Come scrisse Nietzsche,  \textit{abbiamo l'arte per non morire di verit\`{a}}. E la memoria \`{e} forse la pi\`{u} antica e raffinata delle arti umane: l'arte di raccontarci una storia abbastanza vera da essere utile, abbastanza falsa da essere sopportabile. C'\`{e} una ragione profonda per cui siamo l'animale dell'oblio. Le api dimenticano i fiori secchi, gli scoiattoli le noci marce, ma noi siamo gli unici consapevoli di dimenticare, gli unici capaci di studiare il proprio oblio. I sistemi digitali, al contrario, non dimenticano mai, e questo \`{e} esattamente il loro limite: la rete conserva tutto allo stesso livello, la foto di ieri vale quanto quella di dieci anni fa, nessuna gerarchia di rilevanza, tutto piatto e congelato in un eterno presente senza profondit\`{a}. Noi siamo vivi proprio perch\'{e} dimentichiamo. Ogni giorno il cervello decide cosa conservare e cosa lasciar svanire, e questa selezione continua \`{e} la vita cognitiva stessa. Forse, un domani, anche le intelligenze artificiali dovranno imparare a dimenticare per diventare davvero intelligenti. Per ora, quella capacit\`{a} resta unicamente nostra: un dono evolutivo travestito da difetto, una benedizione che ci ostiniamo a chiamare errore.

Questa consapevolezza della fallibilit\`{a} mnemonica solleva per\`{o} una domanda inquietante: se non possiamo fidarci dei nostri ricordi, possiamo almeno fidarci del nostro giudizio presente? Una cosa \`{e} ammettere che la memoria ricostruisce il passato, un'altra \`{e} accettare che anche le nostre decisioni attuali, quelle che prendiamo a mente lucida e con tutti i dati a disposizione, siano altrettanto inaffidabili. Eppure manteniamo una curiosa asimmetria: accettiamo, magari a malincuore, che la memoria ci inganni, ma quando si tratta di valutare probabilit\`{a} o calcolare rischi ci comportiamo come se il cervello cambiasse improvvisamente modalit\`{a}, come se lo stesso organo incapace di ricordare con esattezza cosa \`{e} successo ieri potesse oggi processare statistiche con precisione matematica.

La neurobiologia ci dice che si tratta dello stesso cervello, con gli stessi limiti, le stesse euristiche, gli stessi bias sistematici. L'ippocampo che frammenta e ricostruisce i ricordi \`{e} connesso alla corteccia prefrontale che valuta opzioni e conseguenze; l'amigdala che colora emotivamente il passato \`{e} la stessa che reagisce alle scelte future. Non esistono compartimenti stagni nel cervello: \`{e} un sistema integrato in cui memoria, emozione e ragionamento si influenzano di continuo. Questa continuit\`{a} suggerisce una verit\`{a} scomoda: se la memoria \`{e} un narratore inaffidabile, il ragionamento potrebbe essere un matematico incompetente. Ed \`{e} proprio l\`{i} che andremo nel prossimo capitolo.